% ============================================================================
% main.tex — DCAT: A Dynamic Cross-Attention Transformer for Conditional
% Multimodal Email Spam Detection
%
% Springer-compatible journal-style manuscript (native layout).
% Class rationale: the Springer Nature sn-jnl class was not retrievable from
% the official template sources at production time (see README.md), so this
% manuscript uses a stable single-column journal layout with standard
% packages, matching the reference papers' geometry (155 x 235 mm trim,
% single column, numbered sections/equations/citations).
%
% Content source (frozen): ../PAPER_CONTENT_REGENERATED.md
% Bibliography: references.bib (33 verified entries, unsrtnat = order of
% first appearance).
% Figures 1-4: native TikZ/pgfplots vector code inline in this file.
% ============================================================================
\documentclass[10pt,twoside]{article}

% ------------------------------- geometry ----------------------------------
% Reference papers (Multimedia Tools and Applications): 155 x 235 mm trim.
\usepackage[papersize={155mm,235mm},top=17mm,bottom=19mm,inner=17.5mm,outer=16.5mm,headsep=4mm,headheight=13pt]{geometry}

% ------------------------------- packages ----------------------------------
\usepackage[utf8]{inputenc}
\usepackage[T1]{fontenc}
\usepackage{amsmath,amssymb,amsfonts}
\usepackage{bm}
\usepackage{graphicx}
\usepackage{booktabs}
\usepackage{multirow}
\usepackage{adjustbox}
\usepackage{caption}
\usepackage{subcaption}
\usepackage{tikz}
\usetikzlibrary{arrows.meta,positioning,fit,calc,backgrounds,matrix}
\usepackage{pgfplots}
\pgfplotsset{compat=1.9}
\usepackage{microtype}
\usepackage{fancyhdr}
\usepackage[numbers,sort&compress]{natbib}
\usepackage{enumitem}
\usepackage{xcolor}
\usepackage[hidelinks,bookmarksnumbered]{hyperref}

\captionsetup{font=small,labelfont=bf,skip=4pt}
\captionsetup[table]{position=above}

% ------------------------------- typography --------------------------------
% Springer-like serif text. newtx is used when available; plain Computer
% Modern fallback keeps the document compilable on minimal TeX installs.
\IfFileExists{newtxtext.sty}{\usepackage{newtxtext}\usepackage{newtxmath}}{\usepackage{lmodern}}

\setlength{\parskip}{1.5pt}
\setlength{\bibsep}{2.5pt}
\setlist{topsep=3pt,itemsep=1.5pt,parsep=0pt}

% ------------------------------- headers -----------------------------------
\pagestyle{fancy}
\fancyhf{}
\fancyhead[L]{\footnotesize\itshape DCAT: Dynamic Cross-Attention Transformer for Multimodal Email Spam Detection}
\fancyhead[R]{\footnotesize\thepage}
\renewcommand{\headrulewidth}{0.4pt}

% ------------------------------- helpers -----------------------------------
\newcommand{\tblsource}[1]{\par\vspace{2pt}{\footnotesize\itshape Source: #1}}
\newcommand{\code}[1]{\texttt{#1}}

% figure palette (defined globally; used across TikZ scopes)
\definecolor{cellhi}{RGB}{33,84,158}
\definecolor{celllo}{RGB}{214,227,242}


% --- tolerance for long artifact paths ----------------------------
\emergencystretch=1.5em

\title{\bfseries DCAT: A Dynamic Cross-Attention Transformer for Conditional Multimodal Email Spam Detection}
\author{Preetham S\\{\normalsize\itshape (Guide: Dr.\ Prathima Mabel)}}
\date{}


\begin{document}
\maketitle
\thispagestyle{fancy}
% ------------------------------- abstract ----------------------------------
\begin{abstract}
\noindent
Email spam remains one of the most persistent security problems in electronic communication, and a growing share of unwanted mail carries its persuasive or malicious payload in images as well as text. Multimodal spam detectors therefore pair a text encoder with an image encoder, but two structural properties of real email corpora are largely unaddressed: the overwhelming majority of messages are text-only, and genuinely multimodal messages are extremely scarce. This paper presents DCAT (Dynamic Cross-Attention Transformer), a conditional multimodal architecture for email spam detection that combines bidirectional cross-attention between projected text and image representations with deterministic modality routing and an optional early-exit classifier that avoids unnecessary vision computation. On the full test set, Full-DCAT achieves a three-seed mean Macro-F1 of 0.9862, compared with 0.9869 for a reimplementation of a published multimodal baseline and 0.9871 for Late Fusion; the differences are small, and the usable corpus is 87.1\% text-only. On the 465-message multimodal test subset, Full-DCAT does not outperform Late Fusion (McNemar p = 1.0000), while DCAT-No-XAttn reaches Macro-F1 0.8317 versus 0.4968 for Full-DCAT; neither difference is statistically significant at $\alpha$ = 0.05 given the subset's size. Five-fold cross-validation yields a mean multimodal Macro-F1 of 0.7077 for Full-DCAT and 0.5453 for Late Fusion, with higher variance for Full-DCAT (0.2639 versus 0.1222). Weighted-DCAT raises multimodal Macro-F1 from 0.6253 to 0.7478, with pooled out-of-fold McNemar tests of p < 10$^{-44}$ against unweighted Full-DCAT and p < 10$^{-5}$ against Late Fusion. The early-exit mechanism skips vision computation for 4.45\% of multimodal emails and reduces average FLOPs by 0.63\% on the evaluated workload. The results indicate that architectural expressiveness alone does not translate into gains when paired training examples are extremely scarce and the corpus is dominated by text-only messages, and that the training objective can matter more than the fusion mechanism.
\end{abstract}

\vspace{2pt}
\noindent{\small  multimodal spam detection \textperiodcentered{} cross-attention \textperiodcentered{} email classification \textperiodcentered{} conditional computation \textperiodcentered{} modality routing \textperiodcentered{} DistilBERT \textperiodcentered{} MobileViT \textperiodcentered{} multimodal fusion \textperiodcentered{} deep learning}

% ============================ body ============================
\section{Introduction}\label{sec:1}

Email is the foundation of asynchronous digital communication. Unsolicited bulk email -- spam -- has been a constant across email history and remains one of the most intractable cybersecurity problems faced by people and organizations \citep{janez2023review}.

Spam is never static: campaigns update their messaging to work around filters; a growing body of literature documents how images -- banners, phishing screenshots, logos, text-as-image -- carry part of the spam signal that text-only filters cannot see \citep{annadatha2018image}, \citep{kyaw2024phishing}.

When the content that distinguishes spam from benign mail is partly in an image, a text-only detector is structurally blind to it.

Multimodal spam detection creates conditions that standard text classifiers do not see. A real-world email corpus is heterogeneous in a particular way: most messages are text-only, only some pair text with an image, and far fewer convey their signal mainly in an image. This heterogeneity creates two intertwined problems.

First, a model that encodes every email in a visual backbone makes needless visual computation for the majority that contain only text.

Second, a model that learns interactions \emph{between} modalities needs enough paired examples to reliably measure them in natural corpora.

Previous multimodal spam systems typically encode text and images separately then fuse them through concatenation or other late-fusion strategies \citep{zhang2023mmtd}, with attention- and gating-based techniques developed in related multimodal literature \citep{lu2019vilbert}. These models differ in where and what kinds of information they share between modalities, but they rely on two assumptions -- that every message needs both encoders, and that enough paired examples exist to learn a fusion with complex interactions -- that natural email does not support.

This paper presents the Dynamic Cross-Attention Transformer (DCAT), a conditionally multimodal email spam detector designed around the realities of email. DCAT is built on three principles: (1) \emph{bidirectional cross-attention}: where representations are projected into a shared latent space, each modality attends to the other before pooling so that a text feature's representation can depend on the image and vice versa, (2) \emph{deterministic routing}: which assigns each email to one of four categories based on extracted content, and (3) an \emph{early exit}: in which a text-only auxiliary classifier is evaluated first on the message's text, and the full multimodal pathway is executed only when that classifier is uncertain. Deterministic routing ensures the visual pathway is only used when needed, and early exit reduces the multimodal computation required.

The study aims to answer the following questions:

Q1 Is the proposed conditional architecture more accurate than a text-only baseline?

Q2 How does the full model compare to an early-exit model? Q3 Is the performance difference maintained when evaluated on the full corpus?

Q4 Is the class-weighted training objective superior to an unweighted one?

Q5 Is the model able to learn effective interactions between the modalities?

Q6 Can the model operate efficiently in terms of parameters, FLOPs, and latency?

Q7 What is the model's attitude towards information sharing, and how is it affected by the size of the training set?

Q8 Do we observe the vision-skip rate?

Q9 Does the model demonstrate learnability and stability?

The answers to these questions are in the following sections.

\section{Background}\label{sec:2}

This section presents the background principles on which the rest of this paper is built, with each subsection concluding at the point where methods or evaluation take up the thread - no unrelated background material is presented that is not used in the paper.

\subsection{Email as a Multimodal Communication Medium}\label{sec:2.1}

The classifier built in this paper makes predictions over message content - which portions of an email are text, which are images - and this content is made well-defined by the formal structure of the email standard.

The MIME (Multipurpose Internet Mail Extensions) standard \citep{rfc2045} designates email messages as structured objects: a set of headers and a body that can contain multiple typed parts.

Any message can then carry simultaneously a plain text body, an HTML presentation of that body, inline images, and attached files - each a discrete media type, encoded in different parts. Pieces of the MIME tree are traversed by parsing tools and assembled into text and image content. This structure is what makes email inherently multimodal; it also implies that any multimodal email classifier must first perform deterministic content extraction, deciding which parts of each message constitute meaningful text and meaningful images. Section 6.1 explains the extraction pipeline used here; the corpus preserves the original MIME structure \citep{spamassassin2006}.

Three message types make up the corpus and occupy all subsequent discussion: text-only, meaning that there is text and no meaningful image; image-only, meaning that there is a meaningful image and no text; and multimodal, meaning that the message has both (Section 2.2).

Different corpora contain very different proportions of these three message types, with immediate consequences for both computation (Section 2.6) and learning (Section 3.3).

\subsection{Spam, Phishing, and the Modern Threatscape}\label{sec:2.2}

Spam is defined as unsolicited bulk email: messages sent in quantity for marketing or attack-related purposes without the user's prior consent \citep{janez2023review}. Phishing is a fraudulent message designed to imitate a trusted entity in order to steal login credentials or induce other fraudulent action \citep{kyaw2024phishing}, which typically overlaps strongly with spam; modern detection perspectives treat spam and phishing on a single continuum \citep{kyaw2024phishing}, and the property of interest is that both rely strongly on the visual channel: glossy graphics and brand logos in marketing spam, branded logos and login screens in scam email \citep{annadatha2018image}.

Any text rendered in an image cannot be filtered by plain text filters, and any visual branding is a way to carry an impersonation claim that plain text can't. As such, spam and phishing cannot be fully defined by text alone. A second property of the spam domain motivates the evaluation methodology used in Sections 6 and 7: dataset shift.

Spam corpora age quickly because spammer strategies evolve in response to filters, so a model's performance is tied to the distribution of the corpus it was evaluated on, and findings may not transfer across corpora or time periods \citep{janez2023review}.

This property motivates a transparent, multi-protocol evaluation strategy.

\subsection{Content Representations for Text}\label{sec:2.3}

This subsection traces through the evolution of how text has come to be represented by the encoder here; Section 2.4 does the same for images. The representation of text for machine learning has gone through several steps: bag-of-words and TF-IDF vectors made up shallow lexical features that could be consumed by sparse linear or probabilistic classifiers \citep{janez2023review}; word embeddings provided dense, learned continuous representations of symbolic lexical units, capturing the distributional properties of words in context \citep{mikolov2013word}; convolutional neural networks enabled learned features from the local concatenation of characters \citep{zhang2015charcnn} and words \citep{kim2014cnn}; and recurrent neural networks \citep{hochreiter1997lstm} model longer-range sequential dependencies, notably LSTMs. The current generation \citep{devlin2019bert,liu2019roberta,sanh2019distilbert} is based on the Transformer \citep{vaswani2017attention}, which removes the recurrence altogether in favor of multiple stacked layers of selfattention - where every token attends to every other - with learned query, key, and value linear projections, as described in the next subsection.

Pretrained language models (e.g. \citep{devlin2019bert,liu2019roberta}) learn deep, bidirectional contextual representations of large unlabeled corpora that transfer to other NLP problems.

DistilBERT \citep{sanh2019distilbert} is an efficient distilled version of BERT with most of its capabilities but fewer parameters and inference cost, making it appropriate as a shared frozen representation of a text for this work. As of Section 5.2, it is used in the form of its six-layer \code{base-uncased} variant, with 768-dimensional context vectors.

\subsection{Content Representations for Images}\label{sec:2.4}

Visual representations have followed a similar progression: VGG's \citep{simonyan2015vgg} uniform small convolutions demonstrated the benefit of depth, ResNet \citep{he2016resnet} introduced residual pathways to enable even deeper models, and MobileNets \citep{howard2017mobilenets} showed that accuracy suitable for many tasks can be achieved at a mobile-friendly computational budget.

CNNs also form the basis of visual spam detection, which establishes the fact that visual patterns of texture, color, and overlays carry the spam signal of interest. These are recapitulated here for the case of images: \citep{annadatha2018image}. Transformers for images (ViT, \citep{dosovitskiy2021vit}) represented a new point: when treated as an image patch token sequence, and pretrained at scale, a Transformer can be directly trained to produce a classifier.

MobileViT \citep{mehta2022mobilevit} combined the two ideas: a local processing convolution and a global processing Transformer, yielding an efficient universal image encoder.

The current work uses the MobileViT-small \citep{mehta2022mobilevit}, which is available through the \code{timm} PyTorch Image Models library \citep{wightman2019timm}, accompanied by a matching image preprocess pipeline (resize, normalization, etc.). As with text, Section 5.3 describes the model used.

\subsection{Attention and Cross-Modal Attention}\label{sec:2.5}

An attention mechanism allows one vector set Q to produce a weighted combination of another vector set V, where the weights depend on the content of a third set K: each query compares to each key, the relationships are normalized in a softmax, and then each value is proportionally combined using these normalized weights.

In the canonical case \citep{vaswani2017attention}, scaled dotproduct attention:

\begin{equation}\label{eq:1}
\text{Attention}(Q, K, V) = \text{Softmax}\!\left( \frac{QK^{T}}{\sqrt{d_k}} \right) V,
\end{equation}

where $d_k$ is the per-head key dimension. Multi-headed versions (see \citep{vaswani2017attention}) run multiple such attention processes with different linear projections, and concatenate their output. This idea first demonstrated usefulness in sequence-to-sequence translation, where it learned to produce soft alignments between source and target language \citep{bahdanau2015align}, and in vision, where it was used to focus the gaze of image captioners on salient regions \citep{xu2015show}. \emph{Self-attention} on one sequence: this is how the encoders in Sections 2.3 and 2.4 read each individual sequence. \emph{Cross-attention} on two sequences: each element of one sequence learns an attention-weighted sum over the other sequence, conditional on the other sequence; that is, every element in sequence A learns to pay attention to sequence B, and vice versa. When the results of these two directions are learned and combined (A attending to B and B attending to A) the attention flow is \emph{bidirectional}.

Flow of attention in two directions was proposed for machine comprehension to learn query-aware representations of context without treating the query separately and causing an early summarization step \citep{seo2017bidaf}, and the vision-language literature by extending this idea across modalities: split the image and text into separate streams but learn how to pay attention across them via co-attentional Transformer layers \citep{lu2019vilbert}, learn across-modality encoder representations \citep{tan2019lxmert}, or contrastively align modalities instead of applying token-level attention \citep{radford2021clip}.

A layer of bidirectional cross-attention learns representations of each modality that are a function of content from the other modality - this is what property DCAT is taking advantage of in Section 5.5. This is different from \emph{late fusion} approaches where each modality is processed independently until a pooled feature vector is extracted and then fused, or from \emph{gating-based} approaches where learned multiplicative gates determine how a common set of features influences each modality \citep{arevalo2017gmu}.

\subsection{Conditional Computation and Early Exits}\label{sec:2.6}

Standard neural networks apply the same computation for each input. Conditional computation adapts it to the input \citep{han2022dynamicsurvey}, e.g., the composition, depth or width of the network varies depending on the input. Several techniques have been used for it, such as early exiting: placing classifiers at some layer in the network, and if an input can be classified before the last layer, it can exit without completing the forward pass \citep{teerapittayanon2016branchynet}; or dynamic routing: an optimization policy is learned and, for each input, it activates the blocks needed to classify the input (e.g., in the dynamically routed convolutional networks of Skip Net \citep{wang2018skipnet}).

Email classification is a special case with a very simple and natural type of interpretable conditional computation. The -structure of each message is prespecified to be completely known without any neural computation: no vision encoder is used for text-only email, and vice versa. Routing on this deterministic signal is free to learn and avoids all misrouting.

For the broad multimodal market, the full power of a model is available: a text-only auxiliary classifier can be computed first, and the costly vision pathway invoked only when the evidence from the text alone falls below a threshold - an early exit (a ``skip'' that avoids a full encoder rather than a handful of layers).

In this richer case, evaluating the auxiliary classifier and then branching off is additional work, and any realized efficiency gains depend on the workload; the efficiency results in Section 7.7 quantify this.

\section{Related Work}\label{sec:3}

Background and related work answer different questions: what mechanisms does this paper build on?

What mechanisms have other systems built on? What is the specific question this paper asks?

Work is grouped into approach families; no single system anchors this paper.

\subsection{Text-Based Spam Detection}\label{sec:3.1}

Scored pattern rules \citep{spamassassin2006} led and remain the dominant spam defense on which all modern work rests. Supervised traditional classifiers - Naive Bayes, support vector machines, and random forests - over bag-of-words and tf-idf features form the second generation. Character-level \citep{zhang2015charcnn} and word-level \citep{kim2014cnn} learned convolution classifiers and sequential models \citep{hochreiter1997lstm} brought in deep learning. Context-aware transformers - BERT \citep{devlin2019bert}, RoBERTa \citep{liu2019roberta}, and this paper's own minimal DistilBERT \citep{sanh2019distilbert} - comprise the state of the art for text-only email classification, and outperform multimodal systems on text-only corpora (verified in this paper's own evaluation, Section 7.1).

\subsection{Image-Based Spam Detection}\label{sec:3.2}

Image-spam detection treats the image as the primary source of evidence. Convolutional neural networks - from colors and shapes to text overlays - learn spam visual patterns directly from the pixels \citep{annadatha2018image}. Image-identification CNNs - VGG \citep{simonyan2015vgg}, ResNet \citep{he2016resnet}, and MobileNet \citep{howard2017mobilenets} - serve as the feature backbone in such detectors. Image-only classification, however, remains a weak primary signal: images are ambiguous without their associated text, legitimate mail frequently shares similar visual components such as logos and marketing material \citep{janez2023review}, \citep{annadatha2018image}, and this paper's Image-Only variant demonstrates that gap in Section 7.1.

\subsection{Multimodal Spam Detection and Fusion}\label{sec:3.3}

Multimodal spam detectors fuse the two sources of evidence. Encoding the image and text individually and fusing at the end is a proven paradigm: multilingual multimodal spam detector (MMTD) combines a text encoder with an image encoder over text and document images before fusion \citep{zhang2023mmtd}. Interaction is introduced into the fusion mechanism using attention \citep{arevalo2017gmu} and gating: learned gating of the two modalities produces a shared space that learns how the two modalities influence each other, and vision-language models (Section 2.5) - ViLBERT's co-attention \citep{lu2019vilbert}, Cross-modality LXMERT's structured cross-modal encoder \citep{tan2019lxmert}, and CLIP's contrastive vision-language alignment \citep{radford2021clip} - learned token-level interaction across modalities.

However, two features specific to the email corpus divide multimodal detection in the vision-language domain from other situations. First, truly native, natural multimodal email is rare; test set multimodal examples may be too small to estimate saturated modality interaction reliably. Second, the computation pipeline executes the vision encoder on every message, regardless of need. This latter feature is not specific to published systems: the corpus structure that previous spam models themselves articulated \citep{janez2023review}, \citep{annadatha2018image}, \citep{zhang2023mmtd} informs both properties.

\subsection{Conditional Computation in Classification Systems}\label{sec:3.4}

Early-exit architectures introduce additional exits and allow confidently classified samples to skip subsequent computation stages; dynamic-routing architectures learn conditional execution policies at the level of computational blocks; and Han et al. \citep{han2022dynamicsurvey} define these sample-specific execution architectures as a distinct class. The information structure of email, which explicitly identifies the kind of content expected in a message - in this case, the modality or modality set, though this deterministic indication existed in the data collection procedure in Han et al. \citep{han2022dynamicsurvey}'s work - in each message offers a strong, zero-cost signal for content-based modality routing that existing general mechanisms do not consider: the text and image encoders needed for a particular message are known in the message itself, before execution. This paper's approach - using content-conditional modality routing and an early message-level exit to skip vision - is a modality-awareness flavor of conditional computation that is, to our knowledge, the first to appear in email spam detection.

\subsection{Evaluation Practice in Modality-Imbalanced Settings}\label{sec:3.5}

Modality imbalance between available sources of evidence complicates evaluation practice: a model can post excellent results across the full test corpus while failing on the instances that define its unique knowledge, and corpus shift issues common in spam detection mean that common metrics are driven by corpus composition \citep{janez2023review}. These issues motivate the evaluation approach used here: separating full test from multimodal subset performance, performing five-fold cross-validation in the multimodal subset, and testing model paired differences with McNemar \citep{mcnemar1947note} on pooled out-of-sample predictions. \#\#\# 3.6 Summary of Gaps

Four specific gaps emerge from Sections 3.1--3.5, and the paper is organized to address each:

\begin{enumerate}[leftmargin=*,itemsep=1.5pt]
\item[1.] Cross-modal interaction in email spam studies is limited to a few specialized formalisms (if any) \citep{zhang2023mmtd}, \citep{arevalo2017gmu} - none is the dynamic, token-level, bidirectional cross-attention proven in vision-language architectures \citep{lu2019vilbert}, which has not appeared in the surveyed email-spam literature.
\item[2.] Unconditional visual encoding is expensive: multimodal pipelines encode both modalities for every message \citep{zhang2023mmtd} even though real corpora are text-dominated \citep{janez2023review}, \citep{spamassassin2006} --- paying vision cost where it cannot matter.
\end{enumerate}

\begin{enumerate}[leftmargin=*,itemsep=1.5pt]
\item[3.] Class imbalance in the multimodal subset is under-treated. Multimodal spam evaluations do not report imbalance-corrected training or subset-level significance, although the ham/spam imbalance is documented in the domain literature \citep{janez2023review}, \citep{annadatha2018image}.
\item[4.] The gamut of study metrics are rarely paired: report cross-validation stability, report baseline scores, report message and subset test metrics, measure statistical significance - in the cited studies, only some of these were present and when they were, they appeared scattered in the text. DCAT addresses each of these gaps: bidirectional cross-attention (gap 1), deterministic modality routing with early exit (gap 2), a class-weighting ablation (gap 3), and a multi-protocol evaluation with statistical testing (gap 4).
\end{enumerate}

\section{Problem Formulation}\label{sec:4}

\subsection{Task Definition}\label{sec:4.1}

Suppose an email message m has text content t (the email body, subject line, and any other text data that has been extracted) and is accompanied by one or more images $\{I_1, I_2, \ldots, I_K\}$ embedded inline or as attachments. The goal is binary classification of the email message: predict.

\begin{equation}\label{eq:2}
y \in \{0, 1\}, \qquad y = 0 \text{ (ham, legitimate email)}, \quad y = 1 \text{ (spam)}.
\end{equation}

\subsection{Modality Routing}\label{sec:4.2}

Using the features extracted at preprocessing, all emails are routed in a deterministic fashion to one of four routes. Suppose the Boolean indicators $has\_text$ and $has\_image$ denote that meaningful text (10 or more characters in length after cleaning) and at least 1 meaningful image (qualifies the size/quality constraints outlined in section 6.1) are included. Then the rule for routing is:

\begin{equation}\label{eq:3}
\text{modality}(m) = \begin{cases} \text{MULTIMODAL}, & \text{if } has\_text \wedge has\_image, \\ \text{TEXT}, & \text{if } has\_text \wedge \neg has\_image, \\ \text{IMAGE}, & \text{if } \neg has\_text \wedge has\_image, \\ \text{EMPTY}, & \text{otherwise.} \end{cases}
\end{equation}

EMPTY messages - those with neither text nor meaningful image - do not proceed to training or evaluation. The router decision is deterministic: there is no learned classifier, no probability threshold, no learned heuristic. We prefer determinism: the content structure of a message completely decides its modality structure, so a learned router could only add error, cost, and opacity to a decision that content structure already resolves.

\subsection{Conditional Computation Paths}\label{sec:4.3}

DCAT implements three active computation paths:

\begin{itemize}[leftmargin=*,itemsep=1.5pt]
\item \textbf{TEXT path:} text encoder $\rightarrow$ projection $\rightarrow$ mean pooling $\rightarrow$ shared classifier. Do not call the vision encoder.
\item \textbf{IMAGE path:} image encoder $\rightarrow$ projection $\rightarrow$ mean pooling $\rightarrow$ shared classifier. Note: the text encoder is not used.
\item \textbf{MULTIMODAL path:} text encoder $\rightarrow$ projection $\rightarrow$ optional early exit via the text-only classifier $\rightarrow$ if not exited, image encoder $\rightarrow$ projection $\rightarrow$ bidirectional cross-attention $\rightarrow$ mean pooling $\rightarrow$ fusion $\rightarrow$ shared classifier.
\end{itemize}

All trainable parts (projector, cross-attention, classifiers) are the same in each path; only the \emph{order} of operations varies by modality. In Fig. 1, these three paths plus the early-exit branch are shown.

\subsection{Learning Objective}\label{sec:4.4}

The model is trained with the binary cross-entropy over $N$ training examples:

\begin{equation}\label{eq:4}
\mathcal{L} = -\frac{1}{N} \sum_{i=1}^{N} \left[ y_i \log \sigma(\hat{z}_i) + (1 - y_i) \log \left( 1 - \sigma(\hat{z}_i) \right) \right],
\end{equation}

where $\hat{z}_i$ is the predicted logit of email $i$, $\sigma(\cdot)$ is the sigmoid function, and $y_i$ is the label which is 1 if the email $i$ is spam and 0 if it is not. Since the training distribution is ham-heavy (Section 6.2), the weighted variant used by Weighted-DCAT is:

\begin{equation}\label{eq:5}
\mathcal{L}_{weighted} = -\frac{1}{N} \sum_{i=1}^{N} \left[ w_{pos}\, y_i \log \sigma(\hat{z}_i) + (1 - y_i) \log \left( 1 - \sigma(\hat{z}_i) \right) \right],
\end{equation}

with the positive-class weight defined as

\begin{equation}\label{eq:6}
w_{pos} = \frac{N_{ham}}{N_{spam}}.
\end{equation}

This weighting is meant to counteract the ham-heavy class distribution used during training, with its empirical effect observed by the imbalance ablation of Section 7.4.

\section{The DCAT Architecture}\label{sec:5}

\subsection{Architecture Overview}\label{sec:5.1}

DCAT consists of five components: 1. \textbf{Text encoder (DistilBERT \citep{sanh2019distilbert}):} encodes email text into token-level contextual embeddings. 2. \textbf{Vision encoder (MobileViT \citep{mehta2022mobilevit}):} encodes email images into spatial visual feature maps. 3. \textbf{Modality projector:} maps both encoder outputs into a shared latent space of dimension $d_{model} = 256$, using linear projections followed by Layer Normalization. 4. \textbf{Bidirectional cross-attention:} exchanges information between text and image representations in the shared latent space. 5. \textbf{Shared classifier:} maps fused or unimodal shared representations into a spam/ham prediction.

\begin{figure}[t]
\centering
\begin{adjustbox}{max width=\textwidth}
\begin{tikzpicture}[
  font=\scriptsize,
  node distance=4mm and 5mm,
  box/.style={draw, rounded corners=1pt, align=center, inner sep=2.5pt, minimum height=6.5mm},
  frozen/.style={box, fill=blue!6},
  trainable/.style={box, fill=orange!12},
  io/.style={box, fill=gray!10},
  flow/.style={-{Latex[length=1.6mm]}, semithick},
  routeline/.style={flow, dashed},
  lab/.style={font=\tiny\itshape, text=black!60}
]
% input and routing
\node[io] (email) {Input email (MIME)};
\node[io, below=of email] (prep) {Preprocessing \& content extraction};
\node[box, below=of prep, fill=red!8, minimum width=44mm] (router) {Deterministic modality router\\{\tiny (TEXT / IMAGE / MULTIMODAL / EMPTY)}};
\draw[flow] (email) -- (prep);
\draw[flow] (prep) -- (router);

% EMPTY discard
\node[box, below right=3mm and 8mm of router, fill=gray!6, dashed] (empty) {EMPTY: discarded};
\draw[routeline] (router.east) -| (empty.north);

% three path columns
\node[frozen, below left=9mm and 16mm of router.south] (txtenc) {Text encoder\\DistilBERT (frozen)};
\node[frozen, below right=9mm and 16mm of router.south] (imgenc) {Vision encoder\\MobileViT-small (frozen)};

\node[trainable, below=4.5mm of txtenc] (txtproj) {Projection + LayerNorm\\{\tiny (trainable, $d_{model}{=}256$)}};
\node[trainable, below=4.5mm of imgenc] (imgproj) {Projection + LayerNorm\\{\tiny (trainable, $d_{model}{=}256$)}};

\node[box, below=8.5mm of $(txtproj)!0.5!(imgproj)$, fill=orange!12, minimum width=46mm] (xattn) {Bidirectional cross-attention\\{\tiny 2 layers, 4 heads (trainable)}};
\node[box, below=4.5mm of txtproj, fill=gray!8] (pooltxt) {Mean pooling};
\node[box, below=4.5mm of imgproj, fill=gray!8] (poolimg) {Mean pooling};

% early exit
\node[trainable, left=10mm of pooltxt, minimum width=24mm, align=center] (aux) {Auxiliary text-only\\classifier (trainable)};
\node[box, left=7mm of aux, dashed, align=center] (exitgate) {$p_{aux} \geq 0.95$ or\\$p_{aux} \leq 0.05$ ?};

% fusion and classifier
\node[box, below=4.5mm of xattn, fill=gray!8] (avg) {Element-wise average};
\node[box, below=5mm of $(pooltxt)!0.5!(poolimg)$, yshift=-12mm, fill=orange!12, minimum width=40mm] (cls) {Shared classifier (trainable)};
\node[io, below=4.5mm of cls] (out) {Spam / ham};

% flows
\draw[routeline] (router.south west) ++(0,0) |- (txtenc.north);
\draw[routeline] (router.south) -- ++(0,-2mm) -| (imgenc.north);
\draw[routeline] (router.south) -- ++(0,-2.5mm) -| ($(xattn.north)+(0,10mm)$) -- (xattn.north);
\draw[flow] (txtenc) -- (txtproj);
\draw[flow] (imgenc) -- (imgproj);
\draw[flow] (txtproj) -- (pooltxt);
\draw[flow] (imgproj) -- (poolimg);
\draw[flow] (pooltxt.south) -- ++(0,-2mm) -| ($(xattn.north)+(-8mm,4mm)$) -- ($(xattn.north)+(-8mm,0)$);
\draw[flow] (poolimg.south) -- ++(0,-2mm) -| ($(xattn.north)+(8mm,4mm)$) -- ($(xattn.north)+(8mm,0)$);
\draw[flow] (xattn) -- (avg);
\draw[flow] (pooltxt.south) ++(-6mm,-2mm) |- ($(cls.north)+(-9mm,6mm)$) -- ($(cls.north)+(-9mm,0)$);
\draw[flow] (avg) |- ($(cls.north)+(0,3mm)$) -- (cls.north);
\draw[flow] (cls) -- (out);
% early exit flows
\draw[flow] (pooltxt.west) -- (aux.east);
\draw[flow] (aux) -- (exitgate);
\draw[flow, dashed] (exitgate.north) |- node[lab, pos=0.8, above]{exit: classify without vision} ($(cls.north)+(-16mm,0)$);
\draw[flow, dashed] (exitgate.south) to[out=250,in=160] node[lab, pos=0.35, below]{no: run vision pathway} ($(xattn.west)+(0,-1mm)$);
\draw[routeline] (router.west) -| node[lab, pos=0.2, above]{TEXT only} (aux.north);

% legend
\node[box, fill=blue!6, below=2mm of out, minimum width=8mm, inner sep=1.5pt] (leg1) {};
\node[lab, right=1mm of leg1] {frozen};
\node[box, fill=orange!12, right=4mm of leg1, minimum width=8mm, inner sep=1.5pt] (leg2) {};
\node[lab, right=1mm of leg2] {trainable};
\end{tikzpicture}
\end{adjustbox}
\caption{DCAT architecture overview. Deterministic modality routing assigns each email to the TEXT, IMAGE, MULTIMODAL, or EMPTY category; only the required encoders are invoked. For MULTIMODAL emails, a text-only auxiliary classifier may exit early and classify without vision; otherwise bidirectional cross-attention exchanges information between the projected streams before pooling and fusion. Frozen components are shown in blue; trainable components in orange.}
\label{fig:architecture}
\end{figure}

The design rationale follows from Sections 2 and 4. The encoders are frozen pre-trained models, so trainable capacity concentrates in the projection, interaction, and classification components - only 2,008,578 of 73,309,090 parameters (2.74\%) are trainable (Section 5.8). The projector exists because the two encoders emit representations in different spaces (768-dimensional text tokens versus spatial visual feature maps); cross-attention requires a common space in which queries from one modality can be compared against keys from the other. The classifier is \emph{shared} across all three paths so that every path's representation is mapped into the same decision space.

\subsection{Text Encoder}\label{sec:5.2}

The text encoder is DistilBERT-base-uncased \citep{sanh2019distilbert}, a six-layer Transformer with 768-dimensional hidden representations.

Given tokenized input, it produces token-level hidden states:

\begin{equation}\label{eq:7}
\mathbf{H}_{text} = \text{DistilBERT}(\mathbf{x}_{text}).
\end{equation}

where the input is truncated or padded to 256 tokens, and an attention mask indicates padding positions so that masked tokens do not contribute to attention or pooling, and B denotes the batch size.

All encoder parameters are frozen during training: the encoder acts as a fixed feature extractor, and no encoder gradient updates occur.

\subsection{Vision Encoder}\label{sec:5.3}

The vision encoder is MobileViT-small \citep{mehta2022mobilevit}, accessed via the \code{timm} library \citep{wightman2019timm}. Images are resized to $224 \times 224$ pixels before encoding. With normalized image tensor $\mathbf{x}_{img}$, it produces spatial feature representations:

\begin{equation}\label{eq:8}
\mathbf{H}_{img} = \text{MobileViT}(\mathbf{x}_{img}).
\end{equation}

$\mathbf{H}_{img}$ contains $S$ spatial tokens; for the evaluated configuration the feature map corresponds to a $15 \times 15$ spatial grid. Like the text encoder, the vision encoder is frozen during training.

\subsection{Modality Projector}\label{sec:5.4}

The projector maps each modality into the shared $d_{model}$-dimensional latent space ($d_{model} = 256$):

\begin{equation}\label{eq:9}
\mathbf{Z}_{text} = \text{LayerNorm}\!\left( \mathbf{H}_{text} \mathbf{W}_{text} + \mathbf{b}_{text} \right),
\end{equation}

\begin{equation}\label{eq:10}
\mathbf{Z}_{img} = \text{LayerNorm}\!\left( \mathbf{H}_{img} \mathbf{W}_{img} + \mathbf{b}_{img} \right),
\end{equation}

with projection matrices

\begin{equation}\label{eq:11}
\mathbf{W}_{text} \in \mathbb{R}^{768 \times d_{model}}, \qquad \mathbf{W}_{img} \in \mathbb{R}^{D_{img} \times d_{model}},
\end{equation}

where $D_{img}$ is the MobileViT feature dimension and $d_{model} = 256$. The projections are trainable; they are the first stage at which the two modalities inhabit a common space.

Layer normalization after the linear projection stabilizes the scale of both projected streams before they interact, which matters because the two encoders were pre-trained with different internal distributions.

\subsection{Bidirectional Cross-Attention}\label{sec:5.5}

The cross-attention module has a total of $N_{xattn} = 2$ identical layers, each performing information exchange in both directions, sharing the co-attention design of \citep{lu2019vilbert} and the scaled dot-product approach of Eq. (1).

For text-to-image attention, the queries are obtained from the text stream, and the keys and values from the image stream:

\begin{equation}\label{eq:12}
\mathbf{Q}_{text} = \mathbf{Z}_{text} \mathbf{W}_{Q,text}, \qquad \mathbf{K}_{img} = \mathbf{Z}_{img} \mathbf{W}_{K,img}, \qquad \mathbf{V}_{img} = \mathbf{Z}_{img} \mathbf{W}_{V,img},
\end{equation}

\begin{equation}\label{eq:13}
\mathbf{Z}_{img \rightarrow text} = \text{Softmax}\!\left( \frac{\mathbf{Q}_{text} \mathbf{K}_{img}^{T}}{\sqrt{d_k}} \right) \mathbf{V}_{img}.
\end{equation}

The new, content-aware representation of each text token is therefore re-derived as a mixture of image tokens: what parts of an image are relevant for a specific text token are not predetermined, but learned. This time, the attention from image to text works the other way:

\begin{equation}\label{eq:14}
\mathbf{Q}_{img} = \mathbf{Z}_{img} \mathbf{W}_{Q,img}, \qquad \mathbf{K}_{text} = \mathbf{Z}_{text} \mathbf{W}_{K,text}, \qquad \mathbf{V}_{text} = \mathbf{Z}_{text} \mathbf{W}_{V,text},
\end{equation}

\begin{equation}\label{eq:15}
\mathbf{Z}_{text \rightarrow img} = \text{Softmax}\!\left( \frac{\mathbf{Q}_{img} \mathbf{K}_{text}^{T}}{\sqrt{d_k}} \right) \mathbf{V}_{text}.
\end{equation}

The two streams are then added with residual connections:

\begin{equation}\label{eq:16}
\mathbf{Z}_{text}^{out} = \mathbf{Z}_{text} + \text{MLP}\! \left(\mathbf{Z}_{img\rightarrow text}\right),
\end{equation}

\begin{equation}\label{eq:17}
\mathbf{Z}_{img}^{out} = \mathbf{Z}_{img} + \text{MLP}\!\left(\mathbf{Z}_{text \rightarrow img}\right)
\end{equation}

\begin{figure}[t]
\centering
\begin{adjustbox}{max width=\textwidth}
\begin{tikzpicture}[
  font=\scriptsize,
  node distance=4mm and 6mm,
  box/.style={draw, rounded corners=1pt, align=center, inner sep=2.5pt, minimum height=6mm},
  attn/.style={box, fill=orange!12, minimum width=34mm},
  ln/.style={box, fill=blue!6, minimum width=16mm},
  flow/.style={-{Latex[length=1.6mm]}, semithick},
  lab/.style={font=\tiny\itshape, text=black!60}
]
% input streams
\node[ln] (lnT) {LayerNorm (Pre-LN)};
\node[ln, right=46mm of lnT] (lnI) {LayerNorm (Pre-LN)};
\node[box, above=3.5mm of lnT, fill=blue!6] (inT) {Text stream $\mathbf{Z}_{text}$};
\node[box, above=3.5mm of lnI, fill=blue!6] (inI) {Image stream $\mathbf{Z}_{img}$};

% projections
\node[box, below=4mm of lnT] (qT) {$\mathbf{Q}_{text} = \mathbf{Z}_{text}\mathbf{W}_{Q,text}$};
\node[box, below=4mm of lnI] (kV) {$\mathbf{K}_{img},\ \mathbf{V}_{img}$\\{\tiny (from $\mathbf{Z}_{img}$)}};
\node[box, below=12mm of lnT] (kT) {$\mathbf{K}_{text},\ \mathbf{V}_{text}$\\{\tiny (from $\mathbf{Z}_{text}$)}};
\node[box, below=12mm of lnI] (qI) {$\mathbf{Q}_{img} = \mathbf{Z}_{img}\mathbf{W}_{Q,img}$};

% attention blocks
\node[attn, below=25mm of lnT] (aTI) {Cross-attention (text$\rightarrow$image)\\{\tiny $\text{Softmax}(\mathbf{Q}_{text}\mathbf{K}_{img}^{T}/\sqrt{d_k}\,)\mathbf{V}_{img}$}};
\node[attn, below=25mm of lnI] (aIT) {Cross-attention (image$\rightarrow$text)\\{\tiny $\text{Softmax}(\mathbf{Q}_{img}\mathbf{K}_{text}^{T}/\sqrt{d_k}\,)\mathbf{V}_{text}$}};

% residual MLP updates
\node[box, below=4mm of aTI, fill=orange!12] (mTI) {$+$ residual; MLP};
\node[box, below=4mm of aIT, fill=orange!12] (mIT) {$+$ residual; MLP};
\node[box, below=4mm of mTI, fill=blue!6] (outT) {Updated text stream $\mathbf{Z}_{text}^{out}$};
\node[box, below=4mm of mIT, fill=blue!6] (outI) {Updated image stream $\mathbf{Z}_{img}^{out}$};

\draw[flow] (inT) -- (lnT);
\draw[flow] (inI) -- (lnI);
\draw[flow] (lnT) -- (qT);
\draw[flow] (lnI) -- (kV);
\draw[flow] (qT) -- (aTI);
\draw[flow] (kV) -- (aTI);
\draw[flow] (qI) -- (aIT);
\draw[flow] (kT.south) |- ($(aIT.north)+(6mm,8mm)$) -| ($(aIT.north)+(6mm,0)$);
\draw[flow] (lnT.south) ++(-6mm,0) |- ($(kT.west)$);
\draw[flow] (aTI) -- (mTI);
\draw[flow] (aIT) -- (mIT);
\draw[flow] (mTI) -- (outT);
\draw[flow] (mIT) -- (outI);
% cross links: image stream feeds text-direction attention and vice versa
\draw[flow, dashed] (lnI.east) to[out=-25,in=25] node[lab, pos=0.45, right]{K,V to text direction} ($(mIT.east)+(0,9mm)$);
\draw[flow, dashed] (lnT.west) to[out=205,in=155] node[lab, pos=0.4, left]{K,V to image direction} ($(aTI.west)+(0,6mm)$);

% annotation
\node[lab, below=1mm of outT.south east, anchor=north west, align=left] {$\times\, N_{xattn}=2$ layers;\ $h=4$ heads;\ $d_k=64$};
\end{tikzpicture}
\end{adjustbox}
\caption{One bidirectional cross-attention layer. Each stream's queries attend to the other stream's keys and values (Eqs.~\eqref{eq:12}--\eqref{eq:15}); the attended representation is added back to its own stream through a residual MLP update (Eqs.~\eqref{eq:16}--\eqref{eq:17}). Layer Normalization is applied before each sub-layer (Pre-LN). The layer is repeated $N_{xattn} = 2$ times with $h = 4$ heads and per-head dimension $d_k = 64$.}
\label{fig:crossattention}
\end{figure}

Multi-head attention: we use $h = 4$ heads with dimensions $d_{k} = 64$ per head. The layer uses Pre-LN, where Layer Norm is applied before each sub-layer rather than after, which has been shown to stabilize training of the Transformer \citep{xiong2020layernorm}. The residual additions in Eqs. (16)-(17) are used to preserve each stream's original representation so that the cross-attention components focus on enhancement rather than replacement of the unimodal features; this conceptually interacts with the fusion design described below.

\subsection{Pooling and Fusion}\label{sec:5.6}

Sequence level representations are obtained by mean pooling: $\mathbf{z}_{text} = \text{Mean Pool}(\mathbf{Z}_{text})$ and $\mathbf{z}_{img} = \text{Mean Pool}(\mathbf{Z}_{img})$ With mean pooling over the attention mask, variable-length messages (and masked padding) are weighted proportionally.

The multimodal representation is obtained by taking the elementwise mean of the cross-attention outputs:

\begin{equation}\label{eq:18}
\mathbf{z}_{fused} = \frac{\mathbf{z}_{text}^{out} + \mathbf{z}_{img}^{out}}{2}.
\end{equation}

The fused representation is d\_{model} - dimensional, not 2*feature dimension, as late fusion by concatenation would be. Elementwise averaging keeps the multimodal and the unimodal representations in the same space, and this is what allows a single shared classifier to classify on all three paths.

\subsection{Shared Classifier and Early-Exit Mechanism}\label{sec:5.7}

The shared classifier maps the $d_{model}$-dimensional representation to a single logit, implemented the same way on each path:

\begin{equation}\label{eq:19}
p(y = 1 \mid m) = \sigma\! \left( \text{MLP}_{cls}(\mathbf{H}_{fused}) \right),
\end{equation}

where

\begin{itemize}[leftmargin=*,itemsep=1.5pt]
\item $\mathbf{H}_{fused}$ is the fused multimodal representation or the corresponding pooled unimodal representation.
\end{itemize}

In case of MULTIMODAL emails, DCAT optionally checks a text-only auxiliary classifier before calling the vision encoder: the pooled text representation is passed to

\begin{equation}\label{eq:20}
p_{aux}(y = 1 \mid t) = \sigma\! \left( \text{MLP}_{aux}(\mathbf{H}_{unimodal}) \right),
\end{equation}

and the exit rule is

\begin{equation}\label{eq:21}
\text{exit if } p_{aux} \geq 0.95 \;\;(\text{classify spam}) \quad \text{or} \quad p_{aux} \leq 0.05 \;\;(\text{classify ham}).
\end{equation}

For values in between the two thresholds, the entire multimodal path is performed: text $\rightarrow$ vision $\rightarrow$ cross-attention $\rightarrow$ fusion $\rightarrow$ shared classifier. The auxiliary classifier is separate from the main shared classifier and is trained alongside the model, so training the main task implicitly calibrates the exit. The symmetric 0.95/0.05 thresholds force near-certainty before skipping vision; their impact on the workload is measured in Section 7.7, and their fixed value in Section 9.2.

\subsection{Training Procedure}\label{sec:5.8}

The training setup used is:

\begin{itemize}[leftmargin=*,itemsep=1.5pt]
\item AdamW \citep{loshchilov2019adamw} optimizer with learning rate 1e-4, default weight decay.
\item Cosine annealing schedule.
\item Binary cross-entropy with logits loss (Eq. 4), with optional positive-class weighting (Eqs. 5-6) in Weighted-DCAT.
\item Mixed-precision FP16 training using PyTorchAMP.
\item Batch size of 64. Epochs: 1.
\item Encoder parameters: frozen DistilBERT and MobileViT; remaining (projection, cross-attention, shared classifier) unfrozen.
\item Available hardware: a CUDA-capable NVIDIA GPU.
\item Random seed: fixed across Python, NumPy, and PyTorch for deterministic runs.
\end{itemize}

The total number of parameters in the full-DCAT model is 73,309,090, of which 2,008,578 are trainable - 2.74\% of the total. The trainable component consists of the projection layer, the cross-attention layers, and the classifiers; freezing the encoders is what makes single-epoch training on this corpus feasible, and it keeps the trainable component of similar size to the fusion variants it is compared against.

5.9 Inference Pipeline At inference time, the process is:

\begin{itemize}[leftmargin=*,itemsep=1.5pt]
\item The deterministic router assigns the email to either the TEXT, IMAGE, MULTIMODAL, or EMPTY modality (EMPTY is dropped).
\item Only the encoder(s) corresponding to the assigned modality are run.
\item If IMAGE or MULTIMODAL, run the image encoder and perform the visual cross-attention.
\item If MULTIMODAL, compute the auxiliary classifier (Eq. 20).
\item Run the vision encoder (if not already) and perform the text image cross-attention.
\item Compute the shared classifier output.
\item Threshold the sigmoid output at 0.5 for the label.
\end{itemize}

\section{Experimental Methodology}\label{sec:6}

\subsection{Corpus and Preprocessing}\label{sec:6.1}\label{sec:preprocessing}

The primary evaluation corpus used throughout the paper derives from the SpamAssassin public email corpus \citep{spamassassin2006}. The messages are naturally occurring email messages which preserve their original MIME structure (plain text email messages, HTML email messages, inline images, images as attachments). The preservation of MIME enables the deterministic modality extraction of Section 4.2.

The evaluation corpus contains 75,419 collected records, of which 73,206 are usable after filtering; it occupies approximately 16 GB on disk. A supplementary profile of 6,046 SpamAssassin messages (4,150 ham, 1,896 spam) is used to characterize corpus composition (Table 1).

The corpus was selected because it is naturally occurring, preserves the MIME structure needed for deterministic modality extraction, and it has been used in prior spam-detection research.

Preprocessing operates on raw \path{.eml} files through four deterministic stages:

\begin{enumerate}[leftmargin=*,itemsep=1.5pt]
\item[1.] \textbf{MIME parsing.} For each \path{.eml} file, parse with the email package, walk the MIME tree to get plain text, HTML, and images. Use Beautiful Soup and the lxml parser to reduce HTML to visible text; decode image parts to raw bytes; non-fatal parsing issues are logged.
\item[2.] \textbf{Content filtering and quality control.} Texts with less than 10 characters after stripping whitespace are discarded. Images with less than $20 \times 20$ pixels or less than 1024 bytes are removed, with the corresponding dropping reasons preserved. If there are multiple images with sufficient information, the largest one in byte size is preserved. Pixels count is used if a tie exists.
\item[3.] \textbf{Image processing.} Chosen images are resized to $224 \times 224$ pixels through bicubic sampling, transformed into RGB color space, and normalized via the \code{timm} \citep{wightman2019timm} preprocessing pipeline, that aligns to the input distribution of the pre-trained MobileViT encoder.
\item[4.] \textbf{Text tokenization.} Email text is tokenized with the DistilBERT tokenizer under \code{max\_length = 256,\allowbreak truncation = longest\_first,\allowbreak padding = max\_length} to produce the input IDs and attention mask.
\end{enumerate}

\begin{figure}[t]
\centering
\begin{adjustbox}{max width=\textwidth}
\begin{tikzpicture}[
  font=\scriptsize,
  node distance=4mm and 6mm,
  box/.style={draw, rounded corners=1pt, align=center, inner sep=2.5pt, minimum height=6.5mm},
  io/.style={box, fill=gray!10},
  stage/.style={box, fill=blue!6},
  filt/.style={box, fill=red!6},
  flow/.style={-{Latex[length=1.6mm]}, semithick},
  lab/.style={font=\tiny\itshape, text=black!60}
]
\node[io] (eml) {Raw \code{.eml} files};
\node[stage, below=of eml] (mime) {1.\ MIME parsing\\{\tiny (email, BeautifulSoup + lxml)}};
\node[filt, below=of mime, minimum width=52mm] (qc) {2.\ Content filtering \& quality control\\{\tiny text $\geq$ 10 chars;\ image $\geq 20\times20$ px and $\geq$ 1024 B;}\\{\tiny largest image selected (pixel area tie-breaker)}};
\node[stage, below left=6mm and 6mm of qc.south] (imgp) {3.\ Image processing\\{\tiny resize $224\times224$ (bicubic), RGB,}\\{\tiny \code{timm} normalization}};
\node[stage, below right=6mm and 6mm of qc.south] (tok) {4.\ Text tokenization\\{\tiny DistilBERT tokenizer,}\\{\tiny \code{max\_length} = 256, attention mask}};
\node[box, below=10mm of $(imgp)!0.5!(tok)$, fill=red!8, minimum width=44mm] (assign) {Modality assignment\\{\tiny (deterministic; Eq.~\eqref{eq:3})}};
\node[box, below left=5mm and 8mm of assign.south, fill=gray!8] (t) {TEXT};
\node[box, below=5mm of assign, fill=gray!8] (i) {IMAGE};
\node[box, below right=5mm and 8mm of assign.south, fill=gray!8] (mm) {MULTIMODAL};
\node[box, right=13mm of assign, dashed, fill=gray!5] (e) {EMPTY:\\discarded};

\draw[flow] (eml) -- (mime);
\draw[flow] (mime) -- (qc);
\draw[flow] (qc.south) -- ++(0,-1.5mm) -| (imgp.north);
\draw[flow] (qc.south) -- ++(0,-1.5mm) -| (tok.north);
\draw[flow] (imgp.south) |- ($(assign.west)+(0,2mm)$) -- (assign.west);
\draw[flow] (tok.south) |- ($(assign.east)+(-2mm,2mm)$) -- (assign.east);
\draw[flow] (assign.south) -- ++(0,-1.5mm) -| (t.north);
\draw[flow] (assign.south) -- (i.north);
\draw[flow] (assign.south) -- ++(0,-1.5mm) -| (mm.north);
\draw[flow, dashed] (assign.east) -- ($(e.west)+(0,2mm)$);
\end{tikzpicture}
\end{adjustbox}
\caption{Preprocessing and data-flow diagram. Raw \code{.eml} files pass through the four deterministic stages of Section~\ref{sec:preprocessing} (MIME parsing; content filtering with the size and quality thresholds annotated; image processing; text tokenization) before deterministic modality assignment. EMPTY messages are discarded from training and evaluation.}
\label{fig:preprocessing}
\end{figure}

\subsection{Dataset Composition, Splits, and Leakage Control}\label{sec:6.2}

Table 1 shows the structure of the processed corpus.

\begin{table}[t]
\centering
\caption{Composition profile of the supplementary SpamAssassin subset (6,046 emails)}
\label{tab:composition}
\small
\setlength{\tabcolsep}{4pt}
\begin{tabular}{lrr}\toprule
Property & Count & Percentage \\
\midrule
Total emails & 6,046 & 100.0\% \\
Parse success & 6,046 & 100.0\% \\
Ham (legitimate) & 4,150 & 68.6\% \\
Spam & 1,896 & 31.4\% \\
TEXT & 6,007 & 99.4\% \\
MULTIMODAL & 11 & 0.2\% \\
IMAGE & 1 & 0.02\% \\
EMPTY (excluded) & 27 & 0.4\% \\
TEXT ham & 4,146 & 68.6\% \\
TEXT spam & 1,861 & 30.8\% \\
MULTIMODAL ham & 4 & 0.07\% \\
MULTIMODAL spam & 7 & 0.12\% \\
IMAGE spam & 1 & 0.02\% \\
EMPTY spam & 27 & 0.45\% \\
\bottomrule
\end{tabular}

\end{table}

In this supplementary subset, text-only messages constitute 99.4\% of the processed corpus; just 11 of 6,046 emails are multimodal (4 ham, 7 spam), one email is IMAGE, and 27 are EMPTY and are not used for training and testing. The ham/spam split is 4,150/1,896 messages--a 2.19:1 ham:spam ratio--so the corpus is moderately imbalanced rather than approximately balanced. This profile motivates the working hypothesis of the entire experiment--the full-test measure is ultimately determined by text classification, and the multimodal subset is small by construction. The primary evaluation corpus shows the same structural pattern at scale: of 73,206 usable records, 7,867 (10.7\%) are image-bearing, and the ham/spam imbalance is more pronounced (25,213 ham, 47,993 spam).

The primary evaluation corpus is divided 60,335 / 7,542 / 7,542 records (80\% / 10\% / 10\% of collected records) into training, validation, and test sets with a fixed seed of 42; after EMPTY messages are dropped by the deterministic router, the evaluated test set contains 7,351 messages. The split is group-disjoint: connected components of near-duplicate messages (grouped by text and image hashes) are assigned indivisibly to a single split, so no email or its near-duplicates can appear in more than one set. Sender-leakage and campaign-leakage controls are used where sender or campaign identifiers are available. An example of a lineage entry created under a different 10\% primary-test arrangement is presented in Appendix C; throughout the primary tables, the group-disjoint arrangement above is used.

The leakage analysis tests for five kinds of duplication and overlap: exact textual duplicate pairs across splits; exact image duplicate pairs across splits; perceptual image duplicate pairs across splits, using perceptual hashing; Jaccard near-duplicate text pairs; and sender and campaign name overlap. Exact textual and hash-based deduplication was performed on the entire dataset; the perceptual leakage test was performed on a pilot sample of the data due to the high cost; the evidence from this test therefore supports the deduplication checks applied to the entire data, but should not be considered an exhaustive check of all image pairs.

\subsection{Model Variants}\label{sec:6.3}

Eight variants are evaluated, each serving a defined diagnostic role:

\begin{enumerate}[leftmargin=*,itemsep=1.5pt]
\item[1.] \textbf{All-Spam Prior} --- predicts spam for every input; a class-prior sanity floor.
\item[2.] \textbf{Modality Indicator} --- uses only the modality type as the input feature; tests how much signal the modality label alone carries.3. \textbf{Text-Only} --- activates only the text pathway; the text-dominated corpus's reference model.
\item[4.] \textbf{Image-Only} --- activates only the image pathway; quantifies the standalone visual signal.
\item[5.] \textbf{Late Fusion} --- concatenates mean-pooled DistilBERT and MobileViT representations before classification; the standard independent-encoding fusion design.
\item[6.] \textbf{DCAT-No-XAttn} --- removes cross-attention and averages independently pooled text and image representations; isolates the interaction mechanism from the shared-space design.
\item[7.] \textbf{Full-DCAT} --- the complete bidirectional cross-attention architecture (Section 5).
\item[8.] \textbf{DCAT-Early-Exit} --- Full-DCAT with the auxiliary early-exit mechanism enabled.
\end{enumerate}

Furthermore, a published multimodal baseline is featured as a reimplementation of a late-fusion architecture based on DistilBERT for text and ResNet-18 for images \citep{zhang2023mmtd}. As a research alternative, it is approached in this work as another controlled configuration - same corpus, splits, protocol - within our experiments, rather than as a benchmark.

\subsection{Ablation Design}\label{sec:6.4}

Three ablations isolate the paper's research questions, each tied to a gap from Section 3.6:

\begin{itemize}[leftmargin=*,itemsep=1.5pt]
\item \textbf{Cross-attention ablation} (gap 1): compares Late Fusion, DCAT-No-XAttn, and Full-DCAT on the multimodal subset --- a progression from no interaction, to shared-space averaging, to full bidirectional cross-attention.
\item \textbf{Imbalance ablation} (gap 3): compares Full-DCAT trained with unweighted BCE against Weighted-DCAT with $w_{\text{spam}} = N_{\text{ham}} / N_{\text{spam}}$ (Eqs. 5--6).
\item \textbf{Early-exit ablation} (gap 2): compares Full-DCAT and DCAT-Early-Exit using latency, FLOPs, and vision skip rate.
\end{itemize}

\subsection{Evaluation Metrics}\label{sec:6.5}

The evaluation reports Macro-F1 (the mean of ham-F1 and spam-F1, with each given equal weight), Ham F1, Spam F1, Ham PR-AUC (precision-recall area for ham class), ROC-AUC (overall class separation), and Accuracy (fraction of emails correctly classified). In the multimodal case, we re-evaluate the same metrics based solely on messages labelled as MULTIMODAL (MM-Macro-F1 is the multimodal equivalent of Macro-F1). We choose Macro-F1 as the primary measure, since giving equal weight to the class labels means that it is also influenced by the problem of Majority-Class Collapse (defeating the objective to maximize overall accuracy), as we saw in Section 7.2.

\subsection{Statistical Evaluation}\label{sec:6.6}

Paired model comparisons are conducted with McNemar's test \citep{mcnemar1947note} with a continuity correction (combining the counts of discordant predictions between two models). We use two settings: (1) the multimodal test subset (465 samples, of which 6 are ham) and (2) the pooled out-of-fold predictions from the five-fold cross-validation experiment (~7,867 samples). The multimodal test subset contains only 6 ham messages, too few to support statistical claims from the minority class alone; the pooled setting exists precisely to obtain a better-powered comparison.

\subsection{Efficiency Evaluation}\label{sec:6.7}

Efficiency is measured by total and trainable parameter count; an estimate of the total FLOPs per forward pass; the average per-sample latency; P95 latency; the vision skip-rate for early exit; and the FLOP reduction compared to Full-DCAT. The per-sample vs. total-time metrics are shown separately as they are obtained via different measurement procedures.

\subsection{Error-Analysis Protocol}\label{sec:6.8}

Error analysis displays full-test and multimodal confusion matrices for the main variants of the model, for false positives, false negatives, and instances where one model was wrong but the other was right. Errors for each individual sample are saved to \path{artifacts/error_analysis/dcat_errors.csv} and \path{artifacts/error_analysis/late_fusion_errors.csv}. These individual sample error files are the traceability source for the error discussion in Section 7.9.

\section{Results}\label{sec:7}

With that settled in Section 6, what remains is what the measures reveal. The presentation of results in the section follows the formulation of research questions at the beginning of each subsection. Every figure, reference and table reproduced here came from the evaluation artifacts; if the source artifacts showed contradictions, the contradiction is mentioned here.

\subsection{Full-Test-Set Benchmark}\label{sec:7.1}

\emph{Research question: what are the class signals found in the representations, and how do the fusion designs compare when the full test set dominates the metric?}

\begin{table}[t]
\centering
\caption{Primary benchmark — full test set (three-seed means, seeds 42/43/44)}
\label{tab:benchmark}
\small
\setlength{\tabcolsep}{3.4pt}
\begin{tabular}{lrrrrrr}\toprule
Model & Macro-F1 & Ham F1 & Spam F1 & Ham PR-AUC & ROC-AUC & Accuracy \\
\midrule
All-Spam Prior & 0.4000 & 0.0000 & 0.8000 & 0.0000 & 0.5000 & 0.6600 \\
Modality Indicator & 0.5500 & 0.1000 & 0.8500 & 0.1000 & 0.6000 & 0.7000 \\
Text-Only & 0.9827 & 0.9760 & 0.9895 & 0.9968 & 0.9953 & 0.9854 \\
Image-Only & 0.3220 & 0.4901 & 0.1540 & 0.7183 & 0.5400 & 0.3637 \\
Published Baseline & 0.9869 & 0.9818 & 0.9920 & 0.9975 & 0.9969 & 0.9889 \\
Late Fusion & 0.9871 & 0.9820 & 0.9921 & 0.9972 & 0.9965 & 0.9890 \\
DCAT-No-XAttn & 0.9867 & 0.9815 & 0.9919 & 0.9972 & 0.9964 & 0.9887 \\
Full-DCAT & 0.9862 & 0.9808 & 0.9916 & 0.9973 & 0.9967 & 0.9883 \\
DCAT-Early-Exit & 0.9867 & 0.9814 & 0.9919 & 0.9975 & 0.9968 & 0.9887 \\
\bottomrule
\end{tabular}

{\footnotesize\itshape Source: \path{results/tables/final_results.csv} (three-seed means from seeds 42, 43, and 44).\par}

\end{table}

Text-Only achieves Macro-F1 0.9827. Full-DCAT: 0.9862 versus 0.9869 for the Published Baseline, 0.9867 for DCAT-No-XAttn, and 0.9871 for Late Fusion; on the full test set the difference is small.

One final reporting point (reconciled in publication-readiness phase audit): our Late Fusion figure in the row above is the three-seed mean (0.9871), rechecked versus the identified seed-level source table and seed-master.

The original draft report was transcribing 0.9855 for that row (a copy-paste error, with the designated source table), which was inconsistent with the seed-level analysis in Section 7.8. The two numbers are now consistent. At this level of performance, all text models are within thousandths of one another (on the order of 10$^{-3}$).

Image-Only yields Macro-F1 0.3220 -- below the All-Spam Prior of 0.4000. Because only a small proportion of the corpus has meaningful content in the image, this version of image-only training is being trained on far less data than text-only. The sanity priors are working as expected: the All-Spam Prior gives 0.4000 Macro-F1 (based on the class imbalance of 2.19:1), and the Modality Indicator is giving 0.5500, indicating that there is some signal in the modality label itself, but not as much as in content.

\subsection{Multimodal-Subset Results}\label{sec:7.2}

\emph{Research question: what happens on the multimodal test subset?}

\begin{table}[t]
\centering
\caption{Multimodal subset results}
\label{tab:mm}
\small
\begin{tabular}{lrrr}\toprule
Model & Macro-F1 & Ham F1 & Spam F1 \\
\midrule
All-Spam Prior & 0.4968 & 0.0000 & 0.9935 \\
Modality Indicator & 0.4968 & 0.0000 & 0.9935 \\
Text-Only & 0.5584 & 0.1633 & 0.9535 \\
Image-Only & 0.4968 & 0.0000 & 0.9935 \\
Published Baseline & 0.4968 & 0.0000 & 0.9935 \\
Late Fusion & 0.4968 & 0.0000 & 0.9935 \\
DCAT-No-XAttn & 0.8317 & 0.6667 & 0.9967 \\
Full-DCAT & 0.4968 & 0.0000 & 0.9935 \\
DCAT-Early-Exit & 0.4968 & 0.0000 & 0.9935 \\
\bottomrule
\end{tabular}

\end{table}

The multimodal subset has so few samples that Late Fusion, Full-DCAT, Image-Only, and the All-Spam Prior all predict mostly spam, with Macro-F1 0.4968 and Ham F1 0.0000 --- majority-class collapse. Text-Only reaches 0.5584, the only unimodal model with nonzero ham recall (Ham F1 0.1633), while DCAT-No-XAttn is dominant overall, with Macro-F1 0.8317 and Ham F1 0.6667.

One provenance note is needed. Table 3 reproduces the archived multimodal evaluation table (\path{results/tables/main_results_multimodal.csv}, confirmed by \path{results/tables/modality_analysis.csv}). (The source analysis's own synopsis of this subsection gives Text-Only at 0.4968/0.0000/0.9935 (majority-collapsed) and Image-Only at 0.2863/0.4444/0.1282, values that match neither archived table.)

Thus the summary values seem to refer to an alternate (unrecoverable) aggregation/arrangement; the archived table --- the designated source for these figures --- is presented here.

This conflict is disclosed rather than silently reconciled.

One provenance clarification carries over from the source analysis: the All-Spam Prior accuracy cited on this subset is 0.9871. The primary-test artifacts (\path{artifacts/weighted_primary/results.json}) record a multimodal test subset of 465 messages (459 spam, 6 ham), for which the all-spam prediction yields accuracy 459/465 = 0.9871 and Spam-F1 = 918/924 = 0.9935 -- exactly the archived values. The archived subset metrics in Table 3 are therefore consistent with this artifact-backed multimodal test subset. The supplementary 6,046-message profile (11 multimodal: 4 ham, 7 spam) characterizes corpus composition but is not the evaluation population behind Table 3.

\subsection{Cross-Attention Ablation}\label{sec:7.3}

\emph{Question: does bidirectional cross-attention work on multimodal inputs?}

The cross-attention ablation is done between the same three models as the multimodal analysis: No interaction via late fusion (concatenation-based), DCAT-No-XAttn (shared space, averaged representations), and Full-DCAT (bidirectional cross-attention, a transition from no interaction to completely bi-directionally interacting). These three models show all identical majority-collapsed values (Macro-F1 0.4968, Ham F1 0.0000, Spam F1 0.9935) on the multimodal test subset, except for the DCAT-No-XAttn result (Macro-F1 0.8317, Ham F1 0.6667), which was shown in Table 3 and is the only result from an interaction-capable model (though not interaction-dependent) that was not collapsed.

McNemar's test gives p = 1.0000 for Full-DCAT vs. Late Fusion and p = 0.2482 for Full-DCAT vs. DCAT-No-XAttn.

Neither is less than 0.05.

While there are 465 multimodal test samples but only 6 ham messages among them, under this test alone the comparison of Full-DCAT and Late Fusion is underpowered to statistically distinguish, and the advantage of the simpler DCAT-No-XAttn over Full-DCAT is also statistically insignificant.

\subsection{Class-Imbalance Ablation}\label{sec:7.4}

\emph{Research question: does class weighting matter more than architecture?}

\begin{table}[t]
\centering
\caption{Class-imbalance ablation}
\label{tab:imbalance}
\small
\begin{tabular}{lrrrr}\toprule
Model & Macro-F1 & Spam F1 & Ham F1 & MM-Macro-F1 \\
\midrule
Full-DCAT & 0.9862 & 0.9916 & 0.9809 & 0.6253 \\
Weighted Full-DCAT & 0.9880 & 0.9927 & 0.9833 & 0.7478 \\
\bottomrule
\end{tabular}

\end{table}

Weighted-DCAT improves full-test Macro-F1 from 0.9862 to 0.9880 and multimodal Macro-F1 from 0.6253 to 0.7478. Corresponding pooled out-of-fold McNemar comparisons yield p = 2.49 $\times$ 10$^{-45}$ for Weighted-DCAT vs. Full-DCAT, and p = 3.34 $\times$ 10$^{-6}$ for Weighted-DCAT vs. Late Fusion (both far below $\alpha$ = 0.05), on the 7,867 pooled out-of-fold samples -- the corpus's full multimodal population (327 ham, 7,540 spam).

The imbalance ablation is the only comparison in this study with statistically decisive differences.

\subsection{Multimodal Robustness: Five-Fold Cross-Validation}\label{sec:7.5}

\emph{Research question: is the multimodal signal stable?}

Five-fold cross-validation is performed over the corpus's 7,867 multimodal messages, with out-of-fold predictions pooled across folds (Table 5).

\begin{table}[t]
\centering
\caption{Multimodal robustness (five-fold CV over the 7,867 multimodal messages)}
\label{tab:cv}
\small
\begin{tabular}{lrr}\toprule
Model & MM-Macro-F1 (mean) & MM-Macro-F1 (std) \\
\midrule
DCAT-No-XAttn & 0.5171 & 0.0597 \\
Full-DCAT & 0.7077 & 0.2639 \\
Late Fusion & 0.5453 & 0.1222 \\
Published Baseline & 0.5564 & 0.1192 \\
\bottomrule
\end{tabular}

\end{table}

Full-DCAT has the highest reported mean MM-Macro-F1 at 0.7077; its standard deviation of 0.2639 is also the highest among the listed models. The per-fold results (Section 7.6) show large changes between folds, including folds in which Full-DCAT reached 1.0000 and folds in which it was close to the All-Spam Prior level. The small amount of multimodal data in the evaluation protocol limits how strongly these fold-level differences can be generalized.

\subsection{Per-Fold Behavior}\label{sec:7.6}

\emph{Research question: where does the variance come from?}

\begin{table}[t]
\centering
\caption{Per-fold MM-Macro-F1 from five-fold CV}
\label{tab:folds}
\small
\begin{tabular}{rrrr}\toprule
Fold & Late Fusion & DCAT-No-XAttn & Full-DCAT \\
\midrule
1 & 0.4971 & 0.6647 & 0.4971 \\
2 & 0.4971 & 0.6647 & 0.4971 \\
3 & 0.4968 & 0.6644 & 0.4968 \\
4 & 0.7380 & 0.5451 & 1.0000 \\
5 & 0.4975 & 0.6531 & 1.0000 \\
\bottomrule
\end{tabular}

{\footnotesize\itshape \emph{Source: \path{artifacts/weighted_cv/cv_results.csv}.} The fold-level values above are reproduced from the recorded analysis; the archived cross-validation artifacts record the same models' means under a different fold numbering, so the per-fold ordering could not be independently re-derived from the repository.\par}

\end{table}

The fold-level values vary substantially. Full-DCAT reaches 1.0000 in folds 4 and 5, while folds 1--3 remain near the All-Spam Prior level. The variation is consistent with a metric that is highly sensitive to the composition of a small multimodal evaluation set: the same model is perfect in some folds and majority-collapsed in others, depending on which handful of multimodal messages lands in the fold's evaluation split.

\subsection{Efficiency Results}\label{sec:7.7}

\emph{Research question: what does conditional computation actually save?}

\begin{table}[t]
\centering
\caption{Efficiency comparison}
\label{tab:efficiency}
\small
\setlength{\tabcolsep}{3.6pt}
\begin{tabular}{lrrrr}\toprule
Model & Params & Trainable params & FLOPs & Latency/sample (ms) \\
\midrule
Published Baseline & 77,967,426 & 428,034 & ~13B & 1.48 \\
Late Fusion & 71,761,314 & 460,802 & ~13B & 2.65 \\
DCAT-No-XAttn & 72,784,034 & 1,483,522 & ~13B & 2.66 \\
Full-DCAT & 73,309,090 & 2,008,578 & ~13B & 2.70 \\
DCAT-Early-Exit & 73,309,090 & 2,008,578 & ~11B (avg) & 2.68 \\
\bottomrule
\end{tabular}

\end{table}

DCAT has roughly four times as many trainable parameters as Late Fusion --- 2.01M versus 0.46M --- including additional parameters for the cross-attention blocks and classifier. The total FLOPs with these four always-on variants is roughly 13B each, since the frozen encoders make up the bulk of the computation. About 11B FLOPs over the whole workload for DCAT-Early-Exit (the source artifact annotates this cell ``avg, 90\% text-only'') over the whole workload.

In these measurements, the source artifact cell's number is a rough average over the text-only majority of the corpus, which causes it to be not directly comparable with 12.7B in Table 8, since the early exit produces small savings on the tests that trigger the early exit.

That Published Baseline can be read as having lowest latency per sample at 1.48ms, with Full-DCAT and DCAT-Early-Exit measuring 2.70ms and 2.68ms respectively.

\begin{table}[t]
\centering
\caption{Early-exit efficiency details (full test set)}
\label{tab:earlyexit}
\small
\begin{tabular}{lrrr}\toprule
Metric & Full-DCAT & DCAT-Early-Exit & Diff \\
\midrule
Total latency & 59.9 ms & 74.7 ms & +24.82\% \\
P95 latency & 60.3 ms & 113.6 ms & +88.39\% \\
Total FLOPs & 12.8B & 12.7B & -0.63\% \\
Vision skip rate & 0.00\% & 4.45\% & +4.45\% \\
\bottomrule
\end{tabular}

\end{table}

The vision skip rate is 4.45\%, and the corresponding average FLOP reduction is 0.63\%. The total-time measurement shows an increase of 24.82\% for the early-exit variant, while the per-sample measurement in Table 7 changes from 2.70 ms to 2.68 ms. These values come from different measurement procedures and should not be combined into a single latency claim. The measured overhead of evaluating the auxiliary classifier and handling the conditional branch offsets the theoretical computational reduction in the total-time experiment.

\subsection{Seed Variance and Modality-Specific Performance}\label{sec:7.8}

\emph{Research question: how stable are the headline numbers?}

\begin{table}[t]
\centering
\caption{Three-seed Macro-F1 variance}
\label{tab:seeds}
\small
\begin{tabular}{lrrrrr}\toprule
Model & Seed 42 & Seed 43 & Seed 44 & Mean & Std \\
\midrule
Published Baseline & 0.9870 & 0.9865 & 0.9871 & 0.9869 & 0.0003 \\
Late Fusion & 0.9868 & 0.9869 & 0.9875 & 0.9871 & 0.0003 \\
Full-DCAT & 0.9860 & 0.9863 & 0.9864 & 0.9862 & 0.0002 \\
\bottomrule
\end{tabular}

\end{table}

Across the three seeds, the standard deviation is at most 0.0003 for these models. Late Fusion has the highest mean among the three (0.9871), followed by the Published Baseline (0.9869) and Full-DCAT (0.9862); the differences between the means are small.

\begin{table}[t]
\centering
\caption{Modality-specific Macro-F1 (full test set, seed 42)}
\label{tab:modality}
\small
\begin{tabular}{lrr}\toprule
Model & Text-only emails & Multimodal emails \\
\midrule
Late Fusion & 0.9858 & 0.4968 \\
DCAT-No-XAttn & 0.9871 & 0.8317 \\
Full-DCAT & 0.9865 & 0.4968 \\
\bottomrule
\end{tabular}

\end{table}

On the TEXT modality, the reported models have Macro-F1 values of 0.9858 or higher. The IMAGE modality contains only one sample, so any associated metric is highly sensitive to that single prediction. On the MULTIMODAL subset, DCAT-No-XAttn reports 0.8317 while Full-DCAT and Late Fusion report 0.4968 --- the same ordering the multimodal analyses of Sections 7.2--7.4 exhibit, now isolated from the text-dominated remainder of the corpus.

\subsection{Confusion-Matrix Results}\label{sec:7.9}

\emph{Where do remaining errors aggregate?} Figure 4 shows the archived confusion matrix for DCAT in the weighted-primary configuration (panel a), the setting for which the class-weighting analysis of Section 7.4 was conducted, along with a Macro-F1 comparison over all evaluated variants on the multimodal subset (panel b, values from Table 3). In the weighted-primary setting, the DCAT confusion matrix contains 57 false positives and 25 false negatives on the full test set; false positives hinder legitimate mail, false negatives miss spam, and the per-sample error records of Section 6.8 allow them to be directly traced to individual messages. The figure is meant to make the contrast between the behavior of the full corpus and that of the evaluation subset immediately obvious.

\begin{figure}[t]
\centering
\begin{adjustbox}{max width=\textwidth}
\begin{tikzpicture}
% ---- panel (a): DCAT confusion matrix (weighted-primary, full test set) ----
\begin{scope}
% cells: rows = actual (ham, spam), cols = predicted (ham, spam)
\node[draw, fill=cellhi, text=white, minimum width=20mm, minimum height=13mm, font=\small] (c11) at (0,0) {2196};
\node[draw, fill=celllo, minimum width=20mm, minimum height=13mm, font=\small] (c12) at (2.05,0) {57};
\node[draw, fill=celllo, minimum width=20mm, minimum height=13mm, font=\small] (c21) at (0,-1.35) {25};
\node[draw, fill=cellhi, text=white, minimum width=20mm, minimum height=13mm, font=\small] (c22) at (2.05,-1.35) {5073};
\node[font=\scriptsize] at (-1.55,0) {Actual ham};
\node[font=\scriptsize] at (-1.55,-1.35) {Actual spam};
\node[font=\scriptsize] at (0,0.95) {Predicted ham};
\node[font=\scriptsize] at (2.05,0.95) {Predicted spam};
\node[font=\footnotesize\bfseries, anchor=north] at (1.0,-2.15) {(a) DCAT --- weighted-primary, full test set};
\end{scope}
% ---- panel (b): multimodal Macro-F1 comparison (Table 3 values) ----
\begin{scope}[xshift=7.2cm]
\begin{axis}[
  width=6.6cm, height=5.6cm,
  ybar,
  bar width=5.2pt,
  ymin=0, ymax=1.0,
  ylabel={\scriptsize MM-Macro-F1},
  ylabel style={font=\scriptsize},
  ytick={0,0.2,0.4,0.6,0.8,1.0},
  yticklabel style={font=\tiny},
  symbolic x coords={ASP, MI, Text, Image, PB, LF, NoXAttn, DCAT, EE},
  xtick=data,
  xticklabel style={font=\tiny, rotate=60, anchor=east},
  enlarge x limits=0.08,
  tick style={line width=0.3pt},
  axis line style={line width=0.4pt},
]
\addplot[fill=celllo!70!cellhi, draw=cellhi] coordinates {
  (ASP,0.4968) (MI,0.4968) (Text,0.5584) (Image,0.4968)
  (PB,0.4968) (LF,0.4968) (NoXAttn,0.8317) (DCAT,0.4968) (EE,0.4968)
};
\end{axis}
\node[font=\footnotesize\bfseries, anchor=north] at (1.1,-1.55) {(b) Multimodal Macro-F1 by variant};
\end{scope}
\end{tikzpicture}
\end{adjustbox}
\caption{Results visualization. (a)~Confusion matrix of DCAT in the weighted-primary configuration on the full test set ($n = 7{,}351$; 57 false positives, 25 false negatives). (b)~Multimodal Macro-F1 for all evaluated variants on the multimodal subset, reproducing the values of Table~\ref{tab:mm} (ASP: All-Spam Prior; MI: Modality Indicator; Text: Text-Only; Image: Image-Only; PB: Published Baseline; LF: Late Fusion; NoXAttn: DCAT-No-XAttn; DCAT: Full-DCAT; EE: DCAT-Early-Exit). Majority-collapsed variants sit at 0.4968; DCAT-No-XAttn reaches 0.8317.}
\label{fig:results}
\end{figure}

\section{Discussion}\label{sec:8}

This section interprets the evidence of Section 7. Observed results are distinguished from their plausible interpretations; speculation is distinguished from explicit claims.

\subsection{Why the Full-Test Metric Barely Moves}\label{sec:8.1}

The observed result is a 0.0009 Macro-F1 gap between Full-DCAT (0.9862) and Late Fusion (0.9871) on the full test set, with seed-level s.d. of at most 0.0003 (Tables 2 and 9).

Three corpus-level considerations plausibly explain this minimal change.

First, the usable corpus is 87.1\% text-only, which leaves the single cross-attention module that differentiates Full-DCAT from DCAT-No-XAttn to be used on only ~10.7\% of the data; its impact is limited to a small subset of examples. Second, the 327 multimodal ham messages available in the entire corpus are insufficient for stable estimation of cross-modal interactions, so the module's additional trainable parameters may be sensitive to the set of available multimodal training examples. Third, the shared classifier is trained primarily on text-dominated batches, so representations of the multimodal examples must be mapped into a decision space constructed primarily by text-only examples. These factors are consistent with the minimal full-test improvement for more expressive cross-attention architecture; there may be other explanations, but these are consistent with the patterns of results.

Whether the added expressiveness improves overall performance in this context is an empirical question that Section 7.1-7.4 addressed directly, and Section 8.2 interprets.

\subsection{The Multimodal Subset: Signal, Scarcity, and Sensitivity}\label{sec:8.2}

On the multimodal test subset, DCAT-No-XAttn reports Macro-F1 0.8317 compared with 0.4968 for Full-DCAT (Table 3) --- the largest architecture-linked difference in the study --- and it runs in the opposite direction from the design's intent. This large, counter-directional difference was the critical benchmark driving the story of the study. The simpler variant learns pooled independently-vectored text and image representations rather than applying cross-attention.

Two interpretations for this paradox are consistent with the evidence. With only 11 primary multimodal examples, the estimated contribution is highly unstable and sensitive to any prediction errors. And the additional parameters and the consequent complexity of estimating possible interactions in Full-DCAT may have been problematic for estimating those parameters from the limited sample of multimodal training examples, whereas the preserved separate pools of text and image representations in DCAT-No-XAttn seem more robust.

(The statistical result disciplines both readings: the McNemar comparison between Full-DCAT and DCAT-No-XAttn is not significant at $\alpha$ = 0.05 (p = 0.2482).)

What the small evaluation does establish, conclusively, is the ultimate failure mode: every fully trained interaction-capable model that collapsed did so to majority-class prediction (Ham F1 0.0000), exactly the pattern the Macro-F1 headline metric is chosen to expose.

\subsection{Class Weighting Changes More than Architecture Does}\label{sec:8.3}

The imbalance ablation isolates a training-objective effect that changing the class weights in the loss function changes predictions substantially on the pooled evaluation. Weighted-DCAT improves the multimodal Macro-F1 from 0.6253 to 0.7478 and full-test Macro-F1 from 0.9862 to 0.9880; the pooled out-of-fold McNemar comparisons are statistically significant for both Weighted-DCAT versus Full-DCAT (p = 2.49 $\times$ 10$^{-45}$) and Weighted-DCAT versus Late Fusion (p = 3.34 $\times$ 10$^{-6}$).

These evaluations show that changing the class-weighting produces a greater gain in the low-data, strongly-imbalanced multimodal setting than replacing the cross-attention architecture with the simpler, pooling-only design.

This does not establish that class weighting dominates or outweighs architecture, but it does demonstrate a larger measurable difference on the evaluated multimodal task.

\subsection{Efficiency: Theoretical Savings versus Measured Overhead}\label{sec:8.4}

The early-exit mechanism avoids the vision encoder on 4.45\% of the multimodal messages in this evaluation; because only a small fraction of the corpus is multimodal, the resulting average FLOP reduction over the evaluated workload is 0.63\%. The two latency measures tell a conflicting story: on a per-message basis (Table 7), the saved latency (2.70 ms vs.

2.68 ms) is negligible, but for the entire evaluation workload (Table 8), the total latency increased 24.82\%, with an 88.39\% increase in P95 latency. For a definitive statement, the figures cannot be combined; the total-time metric demonstrates that the ``overhead'' of conditional execution, which is required to evaluate an auxiliary classifier and implement the branch, is sufficient to make the overall prediction workload slightly slower than the baseline. (A corpus with a larger proportion of multimodal messages would provide a different workload for the same mechanism; that case was not evaluated here --- the savings scale with multimodal prevalence, not with the mechanism itself.)

\subsection{What DCAT Does and Does Not Establish}\label{sec:8.5}

Read together, the protocols give a complementary view of the architecture, and the claims can accordingly be safely bounded. What the evidence supports: DCAT matches the strong baselines on the full test set, with a small seed-stable gap (Macro-F1 0.9862 versus 0.9869--0.9871), has the highest across-fold average multimodal Macro-F1 (0.7077 versus 0.5453 for Late Fusion) but with greater variation (Section 7.6), class weighting produces the only statistically significant improvement, and the early-exit achieves workload-dependent reductions at some measured overhead cost. What the evidence does not support: a claim that cross-attention improves full-test results in this corpus/protocol combination - the full-test results show no advantage of cross-attention over the published Late Fusion baseline, and the direct-evidence advantage of the simpler variant on the subset is statistically ambiguous. The most honest summary is that, in this corpus, expressiveness in the architecture was not beneficial, while expressiveness in the training objective was.

The usable corpus is 87.1\% text-only, so these conclusions should not be generalized to email corpora with a higher proportion of image-containing messages.

\section{Limitations and Threats to Validity}\label{sec:9}

\subsection{Dataset Limitations}\label{sec:9.1}

\textbf{Extreme multimodal ham scarcity.} Only 327 of the corpus's 7,867 multimodal emails are ham. Scores on this multimodal subset are affected significantly by single predictions, and cannot be assumed to generalize to any other domain. \textbf{Single corpus.} All experiments were conducted on the SpamAssassin corpus. No performance evaluation was conducted on other collections of email. \textbf{No OCR.} The system does not convert text within images into plain text. Spam messages that use images to display the main body of text have all that text stripped away prior to filtering. \textbf{Single-image constraint.} For each multi-image message only the largest, most meaningful image is forwarded to the model; the remaining visual elements are not displayed to it.

\subsection{Architecture Limitations}\label{sec:9.2}

\textbf{Frozen encoders.} DistilBERT and MobileViT are frozen; they were not fine-tuned on the domain data. \textbf{Fusion method.} The proposed multimodal representation is (element-wise averaging) fused after cross-attention; we did not test other fusion methods. \textbf{Fixed early-exit thresholds.} The thresholds were set to 0.95 and 0.05; the effects of varying these were not investigated.

\subsection{Evaluation Limitations}\label{sec:9.3}

\textbf{Small multimodal test set.} The multimodal test subset contains 465 messages, of which only 6 are ham. Multimodal metrics and McNemar comparisons on this subset rest on very few minority-class examples. \textbf{Single training epoch.} Every model is trained for a single epoch, while different training schedules might result in different optimization behavior. \textbf{No adversarial assessment.} The model will not be tested using adversarial spam intentionally designed to challenge it. \textbf{Efficiency measurement uncertainty.} While the per-sample vs. total-time measurements differ, the original source artifact for the total-time comparison is no longer provided. The efficiency table contains two FLOPs figures with different averaging conventions (a workload-average estimate for the early-exit variant, and a marginal measured value for the comparison), and these can no longer be converted into each other from the given artifacts.

\subsection{Threats to Validity}\label{sec:9.4}

\textbf{Internal validity.} Many potential sources of evaluation contamination are mitigated by the deterministic modality routing and duplicate checks; however the single-epoch training regime and extremely limited multimodal data are still critical constraints. \textbf{Construct validity.} Although Macro-F1, class-specific F1, PR-AUC, and ROC-AUC are common classification metrics, the small sample size used for the calculation of multimodal Macro-F1 makes it imprecise. \textbf{External validity.} One corpus of emails is used in the experiments, which is text dominant; our results may vary with a corpus where more emails have images, or follow another trend of spam.

\section{Conclusion and Future Work}\label{sec:10}

This paper presented DCAT, a conditional multimodal email spam detection architecture combining bidirectional cross-attention, deterministic routing, and optional early-exit computation, evaluated on a naturally paired corpus (SpamAssassin) whose defining property is multimodal ham scarcity (327 ham among 7,867 multimodal emails) under a text-only majority (87.1\% of usable records). The evidence supports three conclusions. First, on the full test set, Full-DCAT achieves Macro-F1 0.9862 against 0.9869 for the published-baseline reimplementation and 0.9871 for Late Fusion across three seeds; the differences are small, and the experiments do not demonstrate a consistent full-test advantage from adding cross-attention.

Second, the multimodal results are highly sensitive to distribution, class imbalance, and the multimodal set size: Full-DCAT and the simple versions collapse to majority-class prediction on the multimodal test subset, five-fold cross-validation yields the highest mean of 0.7077 for Full-DCAT versus 0.5453 for Late Fusion but also the highest variance, and class weighting - the training-objective change - produces the only statistically significant improvements in this study (multimodal Macro-F1 0.6253 $\rightarrow$ 0.7478; pooled out-of-fold McNemar p = 2.49 $\times$ 10$^{-45}$ against Full-DCAT and p = 3.34 $\times$ 10$^{-6}$ against Late Fusion).

Third, the conditional computation behaves as theory predicts but as workload demands: the early exit bypasses vision for 4.45\% of multimodal cases and reduces FLOPs by 0.63\%, and total-time measurement shows a 24.82\% latency increase from the execution overhead. Future work directly addresses these limitations: (1) corpus expansion to larger corpora with vastly more multimodal messages including image-based phishing and marketing emails so that training-objective effects such as class-weighting improvements can be tested at scale, (2) OCR to extract text from within images as an additional information source, (3) fine-tuning of DistilBERT and MobileViT encoders with a systematic learning-rate schedule, (4) multi-image handling, to support multiple salient images per email as opposed to the single largest, (5) learned routing capable of assigning importance to the data sources beyond the content rules of deterministic routing, (6) contrastive pretraining of the projection and cross-attention modules on paired text-image vectors prior to classification, and (7) large-scale efficiency comparison of early exit at corpora with greater multimodal prevalence across multiple exit thresholds.

\section*{Data Availability}

The evaluation artifacts are available in the project repository: per-sample error records at \path{artifacts/error_analysis/dcat_errors.csv} and \path{artifacts/error_analysis/late_fusion_errors.csv}; cross-validation results at \path{artifacts/weighted_cv/cv_results.csv}; primary benchmark results at \path{results/tables/final_results.csv}; and the complete results lineage at \path{artifacts/final_results/results_lineage.csv}. The SpamAssassin public corpus is publicly available \citep{spamassassin2006}.

\section*{Declarations}

\textbf{Competing interests.} The author declares no competing interests.

\appendix

\section{Dataset and Preprocessing Reference}\label{app:A}

\begin{table}[h]
\centering
\caption{Dataset and preprocessing reference values}
\label{tab:appdataset}
\small
\begin{tabular}{ll}\toprule
Parameter & Value \\
\midrule
Total emails (SpamAssassin supplementary profile) & 6,046 \\
Ham emails & 4,150 (68.6\%) \\
Spam emails & 1,896 (31.4\%) \\
Text-only emails & 6,007 (99.4\%) \\
Multimodal emails & 11 (0.2\%) \\
Image-only emails & 1 (0.02\%) \\
Empty emails (excluded) & 27 (0.4\%) \\
Text truncation length & 256 tokens \\
Image resolution & 224 $\times$ 224 pixels \\
Train/Val/Test split & 70/15/15 (primary configuration) \\
\bottomrule
\end{tabular}

\end{table}

\section{Model Hyperparameters}\label{app:B}

\begin{table}[h]
\centering
\caption{Model hyperparameters}
\label{tab:apphyper}
\small
\begin{tabular}{ll}\toprule
Hyperparameter & Value \\
\midrule
Text encoder & DistilBERT-base-uncased (frozen) \\
Vision encoder & MobileViT-small via timm (frozen) \\
Shared latent dimension & 256 \\
Cross-attention layers & 2 \\
Cross-attention heads & 4 \\
Cross-attention d\_model & 256 \\
Cross-attention d\_ff & 1024 \\
Cross-attention dropout & 0.1 \\
Cross-attention activation & GELU \\
Cross-attention normalization & Pre-LayerNorm \\
Classifier hidden dim & 128 \\
Classifier dropout & 0.3 \\
Classifier activation & GELU \\
Pooling strategy & Mean pooling \\
Optimizer & AdamW \\
Learning rate & 1 $\times$ 10$^{-4}$ \\
Scheduler & Cosine annealing \\
Batch size & 64 \\
Epochs & 1 \\
Precision & FP16 mixed \\
Loss function & BCEWithLogitsLoss \\
Classification threshold & 0.5 \\
Early-exit high threshold & 0.95 \\
Early-exit low threshold & 0.05 \\
Seeds & 42, 43, 44 \\
\bottomrule
\end{tabular}

\end{table}

\section{Results Lineage and Reproducibility}\label{app:C}

\emph{Source: \path{artifacts/final_results/results_lineage.csv}.}

\begin{table}[h]
\centering
\caption{Results lineage (reproduced as recorded)}
\label{tab:lineage}
\scriptsize
\setlength{\tabcolsep}{3pt}
\resizebox{\textwidth}{!}{%
\begin{tabular}{p{1.7cm}p{2.0cm}p{1.9cm}p{1.1cm}p{3.0cm}rp{2.2cm}}\toprule
Metric & Model & Split & Seed/Fold & Source artifact & Value & Aggregation \\
\midrule
Macro-F1 & Full-DCAT & Primary Test (10\%) & Seed 42 & \path{artifacts/experiments/weighted_primary/results.json} & 0.9868 & Unweighted mean of Spam and Ham F1 \\
MM-Macro-F1 & Full-DCAT & Multimodal CV (5-Fold) & 5-Fold CV & \path{artifacts/experiments/weighted_cv/cv_results.csv} & 0.6609 & Mean across 5 folds \\
MM-Macro-F1 & Late-Fusion & Multimodal CV (5-Fold) & 5-Fold CV & \path{artifacts/experiments/weighted_cv/cv_results.csv} & 0.5564 & Mean across 5 folds \\
MM-Macro-F1 & Weighted-DCAT & Multimodal CV (5-Fold) & 5-Fold CV & \path{artifacts/experiments/weighted_cv/cv_results.csv} & 0.6824 & Mean across 5 folds \\
McNemar p-value & Weighted-DCAT vs. Full-DCAT & Pooled OOF (7,867 samples) & All folds & task-1070 log & 2.49 $\times$ 10$^{-45}$ & Exact/asymptotic McNemar with continuity correction \\
McNemar p-value & Weighted-DCAT vs. Late-Fusion & Pooled OOF (7,867 samples) & All folds & task-1070 log & 3.34 $\times$ 10$^{-6}$ & Exact/asymptotic McNemar with continuity correction \\
\bottomrule
\end{tabular}}

{\footnotesize\itshape \emph{Note.} The lineage table is reproduced as recorded. Its seed-42 primary-test entry (0.9868) and CV means (0.6609 / 0.5564 / 0.6824) differ from the three-seed means of Table 2 and the CV means of Table 6 because the lineage rows document different aggregation settings (a seed-42 primary-test run and an alternative CV configuration); the manuscript's headline tables cite their own authoritative sources as noted beneath each table. The split labels reflect the full raw corpus (approximately 73,500 emails, 10\% $\approx$ 7,351 test emails) rather than the deduplicated 6,046-message evaluation corpus (15\% $\approx$ 907), which explains the label arithmetic. A former discrepancy between the Late Fusion three-seed mean in Table 2 (initially 0.9855) and the seed-level analysis (0.9871) was resolved during the publication-readiness audit as a transcription error against the designated source table; the corrected value is used throughout.\par}

\end{table}

% --- flush appendix floats before the bibliography ------------------
\clearpage
\bibliographystyle{unsrtnat}
\bibliography{references}
\end{document}
